\SourcePage{580}{583}
\section{The shell model}

Many nuclear properties admit a good description in terms of the \emph{shell
model}, whose basic picture resembles the description of an atom's electron
shell. In this approach, each nucleon is treated as moving in the
self-consistent field produced by all the remaining nucleons. Because nuclear
forces have a short range, this field falls off rapidly beyond the volume
enclosed by the nuclear ``surface.'' The state of the entire nucleus is
therefore specified by listing the states occupied by its individual
nucleons.

\SourcePage{581}{584}
The self-consistent field has spherical symmetry, with its center naturally at
the nuclear center of mass. This gives rise, however, to a difficulty. In the
self-consistent-field method, the system's wave function is constructed as a
product of one-particle wave functions (or as a suitably symmetrized sum of
such products). A function of this sort does not keep the center of mass at
rest: although it yields zero for the mean velocity of the center of mass, it
also assigns finite probabilities to nonzero values of that
velocity\footnote{For electrons in an atom this difficulty does not arise at
all, since the center of mass is automatically kept at rest by its coincidence
with the position of the stationary heavy nucleus.}.

The difficulty can be circumvented by first removing the motion of the center
of mass whenever a physical quantity is evaluated with the
self-consistent-field wave functions
$\psi(\mathbf r_1,\ldots,\mathbf r_A)$. Let
$f(\mathbf r_i,\mathbf p_i)$ denote any physical quantity, regarded as a
function of the nucleon coordinates and momenta. To evaluate its matrix
elements with the functions $\psi$, one leaves $\psi(\mathbf r_i)$ unchanged
and instead replaces the arguments of $f$ according to
\begin{equation}
 \mathbf r_i\to\mathbf r_i-\mathbf R,\qquad
 \mathbf p_i\to\mathbf p_i-\frac{\mathbf P}{A},
 \tag{118.1}
\end{equation}
where $\mathbf R$ is the radius vector of the nuclear center of mass, $A$ is
the number of particles in the nucleus, and $\mathbf P$ is the momentum of the
nucleus as a whole. The second replacement in (118.1) amounts to subtracting
the center-of-mass velocity $\mathbf V$ from the nucleon velocities,
$\mathbf v_i\to\mathbf v_i-\mathbf V$; here $\mathbf P$ and $\mathbf V$ are
related by $\mathbf P=Am_p\mathbf V$ (S.~Gartenhaus, C.~Schwartz, 1957).

For example, the nuclear dipole-moment operator is
$\hat{\mathbf d}=e\sum_p\mathbf r_p$, the sum extending over every proton in
the nucleus. In self-consistent-field matrix elements, however, it must be
replaced by $e\sum_p(\mathbf r_p-\mathbf R)$. The coordinate of the nuclear
center is
\[
 \mathbf R=\frac1A\left(\sum_p\mathbf r_p+\sum_n\mathbf r_n\right),
\]
where the two sums run over all protons and all neutrons, respectively. Since
the nucleus contains $Z$ protons, the final replacement of the dipole-moment
operator is
\begin{equation}
 e\sum_p\mathbf r_p\to
 e\left(1-\frac ZA\right)\sum_p\mathbf r_p
 -e\frac ZA\sum_n\mathbf r_n.
 \tag{118.2}
\end{equation}

\SourcePage{582}{585}
Thus the protons occur with an ``effective charge'' $e(1-Z/A)$, while the
neutrons have the ``charge'' $-eZ/A$. As (118.2) shows, the correction terms
that enter the dipole moment have relative order of magnitude unity. By
contrast, it is readily seen that the corrections to the magnetic moment and
to the higher electric multipole moments are of relative order $\sim1/A$.

In the nonrelativistic approximation, a nucleon's interaction with the
self-consistent field is independent of its spin. Such a dependence could
only be represented by a term proportional to $\hat{\mathbf s}\mathbf n$,
where $\mathbf n$ is the unit vector along the nucleon's radius vector
$\mathbf r$; this product is a pseudoscalar rather than a true scalar.

Velocity-dependent relativistic terms nevertheless introduce a spin
dependence into the nucleon energy. The largest such term is linear in the
velocity. The three vectors $\mathbf s$, $\mathbf n$, and $\mathbf v$ can be
combined into the true scalar $(\mathbf n\times\mathbf v)\mathbf s$.
Consequently the operator for a nucleon's \emph{spin--orbit coupling} in the
nucleus has the form
\begin{equation}
 \hat V_{sl}=-\varphi(r)(\mathbf n\times\hat{\mathbf v})\hat{\mathbf s},
 \tag{118.3}
\end{equation}
where $\varphi(r)$ is some function of $r$ (compare also the note on p.~579).
Since $m_p(\mathbf r\times\mathbf v)$ is the particle's orbital angular
momentum $\hbar\hat{\mathbf l}$, (118.3) may also be written
\begin{equation}
 \hat V_{sl}=-f(r)\hat{\mathbf l}\hat{\mathbf s},
 \tag{118.4}
\end{equation}
with $f=\hbar\varphi/(rm_p)$. Notice that this interaction is of first order
in $v/c$, whereas the spin--orbit coupling of an atomic electron is a
second-order effect (\S~72). The reason for this distinction is that nuclear
forces already depend on spin in the nonrelativistic approximation, whereas
the nonrelativistic interaction between electrons (the Coulomb forces) has no
spin dependence.

The energy of the spin--orbit interaction is concentrated chiefly near the
nuclear surface; in other words, $f(r)$ decreases toward the nuclear interior.
Indeed, an interaction of this form could not exist at all in unbounded nuclear
matter. The homogeneity of such a system leaves no preferred direction along
which the vector $\mathbf n$ could point.

The interaction (118.4) splits a nucleon level of orbital angular momentum
$l$ into two levels whose angular momenta are $j=$

\SourcePage{583}{586}
$l\pm1/2$. Since
\begin{equation}
 \begin{aligned}
  \mathbf l\mathbf s&=l/2 &&\text{for }j=l+1/2,\\
  \mathbf l\mathbf s&=-(l+1)/2 &&\text{for }j=l-1/2,
 \end{aligned}
 \tag{118.5}
\end{equation}
by (31.3), the magnitude of the splitting is
\begin{equation}
 \Delta E=E_{l-1/2}-E_{l+1/2}=\overline{f(r)}(l+1/2).
 \tag{118.6}
\end{equation}
Experiment places the $j=l+1/2$ level (parallel $\mathbf l$ and $\mathbf s$)
deeper than the $j=l-1/2$ level. Hence $f(r)>0$.

A nucleon's spin--orbit coupling within a nucleus is weak relative to its
interaction with the self-consistent field. At the same time it is generally
larger than the energy of direct interaction between two nucleons in the
nucleus, because the latter decreases more rapidly as the atomic weight
increases.

This hierarchy of interaction energies requires nuclear levels to be
classified by the $jj$-coupling scheme. For each nucleon its spin and orbital
angular momentum combine to form the total angular momentum
$\mathbf j=\mathbf l+\mathbf s$. This is a definite quantity because the
coupling between $\mathbf l$ and $\mathbf s$ is not destroyed by the direct
interaction of the particles with one another (M.~G\"oppert-Mayer, 1949;
O.~Haxel, J.~H.~Jensen, H.~E.~Suess, 1949)\footnote{Only in the lightest
nuclei is the coupling closer to the $LS$ type.}.
The individual nucleons have angular-momentum vectors $\mathbf j$.
These vectors are then added together.
Their sum is the total nuclear angular momentum $\mathbf J$.
This quantity is usually called simply the \emph{nuclear spin}.
The name treats the nucleus as though it were an elementary particle.
In this respect, nuclear levels and atomic levels are classified quite differently.
Consider the electron shell of an atom.
Its relativistic spin--orbit coupling is generally weak.
One comparison is with direct electric interactions.
These interactions are stronger.
Exchange interactions are likewise stronger.
Atomic levels are therefore usually classified by $LS$ coupling.


Each nucleon state in a nucleus is characterized by its angular momentum $j$
and parity. Although neither $\mathbf l$ nor $\mathbf s$ is separately
conserved, the magnitude of the nucleon's orbital angular momentum is still
definite. Indeed, a given $j$ can arise either from a state with $l=j-1/2$ or
from one with $l=j+1/2$. For fixed half-integral $j$, these two states have
opposite parities $(-1)^l$; specifying $j$ and the parity thus fixes the
quantum number $l$ as well.

\SourcePage{584}{587}
Nucleon states having the same $l$ and $j$ are conventionally numbered, in
order of increasing energy, by a ``principal quantum number'' $n$ that takes
integer values beginning with 1\footnote{This differs from the convention for
atomic electron levels, in which the allowed values of $n$ begin with
$l+1$.}. The various states are denoted by $1s_{1/2}$, $1p_{1/2}$,
$1p_{3/2}$, and so forth. The numeral preceding the letter is the principal
quantum number; the letters $s$, $p$, $d$, \ldots indicate $l$ in the usual
way; and the subscript on the letter gives $j$. A state with specified $n$,
$l$, and $j$ can contain at most $2j+1$ neutrons and the same number of
protons.

For a specified configuration, the state of the nucleus as a whole is
customarily written as a numeral giving $J$, with a $+$ or $-$ subscript to
show the state's parity. In the shell model that parity is determined by
whether the algebraic sum of the $l$ values of all nucleons is even or odd.

Analysis of experimental information about nuclear properties reveals a
number of regularities in the arrangement of nuclear levels.

First, nucleon-level energies rise as the orbital angular momentum $l$
increases. This rule follows because a larger $l$ increases the particle's
centrifugal energy and hence lowers its binding energy.

Second, at fixed $l$, the level with $j=l+1/2$ (corresponding to parallel
$\mathbf l$ and $\mathbf s$) lies deeper than the level with $j=l-1/2$. This
rule was already mentioned above in connection with the properties of a
nucleon's spin--orbit coupling inside a nucleus.

A further rule concerns nuclear isotopic spin. Recall that the projection
$T_\zeta$ is fixed by the weight and number of the nucleus (see (116.1)). For
a specified $T_\zeta$, the isospin magnitude may assume any value satisfying
$T\geqslant|T_\zeta|$. Ordinarily the nuclear ground state has the smallest
allowed isospin, namely
\begin{equation}
 T_{\text{ground}}=|T_\zeta|=(1/2)(N-Z).
 \tag{118.7}
\end{equation}
This rule is tied to the character of the neutron--proton interaction: in an
$np$ system, the $T=0$ state (the deuteron state) has a greater binding energy
than the $T=1$ state (see the note on p.~580).

Some rules can also be stated for the spins of nuclear ground states. They
specify how the angular momenta $\mathbf j$ of the individual nucleons combine
into the total nuclear spin. These rules express the tendency of protons and
neutrons that occupy identical nuclear sta\SourcePageJoin{585}{588}tes to form
$pp$ and $nn$ pairs whose angular momenta are opposite. The binding energy of
such a pair is of order 1--2 MeV.

One consequence is that, if both the proton number and the neutron number are
even (an \emph{even--even nucleus}), all nucleon angular momenta cancel in
pairs and the total nuclear angular momentum vanishes.

If either the number of protons or the number of neutrons is odd, and all the
nucleons outside the filled shells occupy identical states, the total nuclear
angular momentum usually equals that of a single nucleon. It is as though,
after every possible proton and neutron pair had formed, just one nucleon with
uncancelled angular momentum remained. The total angular momentum of any
filled shell is automatically zero.

For \emph{odd--odd nuclei}, in which both $Z$ and $N$ are odd, no sufficiently
general rule determines the ground-state spin.

A detailed discussion of the particular sequence in which nuclear shells are
filled would require extensive analysis of the available experimental data
and lies outside the scope of this book. We shall give only a few further
general observations.

In studying atoms, we found that their electron states can be divided into
groups such that completing one group and beginning the next reduces an
electron's binding energy. Nuclei exhibit an analogous pattern, with nucleon
states distributed among the following groups:
\begin{equation}
 \begin{array}{@{}ll@{}}
  1s_{1/2} & 2\text{ nucleons},\\
  1p_{3/2},\ 1p_{1/2} & 6\text{ nucleons},\\
  1d_{5/2},\ 1d_{3/2},\ 2s_{1/2} & 12\text{ nucleons},\\
  1f_{7/2},\ 2p_{3/2},\ 1f_{5/2},\ 2p_{1/2},\ 1g_{9/2}
    & 30\text{ nucleons},\\
  2d_{5/2},\ 1g_{7/2},\ 1h_{11/2},\ 2d_{3/2},\ 3s_{1/2}
    & 32\text{ nucleons},\\
  2f_{7/2},\ 1h_{9/2},\ 1i_{13/2},\ 2f_{5/2},\ 3p_{3/2},\ 3p_{1/2}
    & 44\text{ nucleons}.
 \end{array}
 \tag{118.8}
\end{equation}
For each group the table gives the total number of available proton or neutron
places. It follows that one of these groups is completed whenever the total
proton number $Z$, or the total neutron number $N$, reaches one of the values
\[
 2,\ 8,\ 20,\ 50,\ 82,\ 126.
\]
These are conventionally called \emph{magic numbers}\footnote{The
$1f_{7/2}$ states, with their 8 places, are sometimes separated into a group
of their own, reflecting the fact that 28 too has, to some extent, the
properties of a magic number.}.

\SourcePage{586}{589}
The so-called \emph{doubly magic nuclei}, for which both $Z$ and $N$ are magic
numbers, possess exceptional stability. Relative to neighboring nuclei, their
affinity for one additional nucleon is anomalously small, while their first
excited levels are anomalously high\footnote{Examples are
${}^{4}_{2}\mathrm{He}_{2}$, ${}^{16}_{8}\mathrm{O}_{8}$,
${}^{40}_{20}\mathrm{Ca}_{20}$, and
${}^{208}_{82}\mathrm{Pb}_{126}$. The ${}^{4}\mathrm{He}$ nucleus cannot
bind an additional nucleon at all. The number at the lower right is the value
of $N$ in the nucleus.}.

Within each group in (118.8), the states have been listed approximately in the
sequence in which they are progressively filled along the series of nuclei.
In reality, substantial irregularities occur in this filling. It must also be
remembered that, in heavy nuclei far from magic numbers, the separations
between different levels may be comparable with the ``pairing energy.'' Under
these conditions, the very notion of individual states for the two members of
a pair largely loses its meaning.

We now make several remarks on calculating a nucleus's magnetic moment in the
shell model. By the magnetic moment of a nucleus we naturally mean the moment
averaged over the motion of the particles within it. The mean magnetic moment
$\overline{\boldsymbol\mu}$ is evidently directed along the nuclear spin
$\mathbf J$, since the latter is the nucleus's only preferred direction. Its
operator is therefore
\begin{equation}
 \hat{\overline{\boldsymbol\mu}}=\mu_0g\hat{\mathbf J},
 \tag{118.9}
\end{equation}
Here $\mu_0$ is the nuclear magneton.
The symbol $g$ denotes the gyromagnetic factor.
Consider the projection of this moment.
Its eigenvalue is $\overline\mu_z=\mu_0gM_J$.
Conventionally, $\mu$ denotes the nuclear magnetic moment.
Compare this usage with (111.1).
It is defined as the largest projection value.
Thus $\mu=\mu_0gJ$.
With this notation, one obtains:

\begin{equation}
 \hat{\overline{\boldsymbol\mu}}=\mu(\hat{\mathbf J}/J).
 \tag{118.10}
\end{equation}

The nuclear magnetic moment is composed of the moments of nucleons outside
filled shells, since those of nucleons in filled shells cancel one another.
Every nucleon contributes a magnetic moment having two parts: a spin part and,
for a proton, an orbital part. It is thus expressed as the sum
$g_s\hat{\mathbf s}+g_l\hat{\mathbf l}$.
(Here and below we omit the factor $\mu_0$. We follow the usual convention.
Magnetic moments are understood to be measured in units of the nuclear magneton.)

The orbital and spin gyromagnetic factors are
$g_l=1$, $g_s=5{,}585$ for a proton, and $g_l=0$, $g_s=-3{,}826$ for a
neutron.

\SourcePage{587}{590}
After averaging over a nucleon's motion in the nucleus, its magnetic moment is
proportional to $\mathbf j$. Writing it as $g_j\hat{\mathbf j}$, we have
\[
 g_j\hat{\mathbf j}=g_s\hat{\mathbf s}+g_l\hat{\mathbf l}
 =\frac12(g_l+g_s)\hat{\mathbf j}
 +\frac12(g_l-g_s)(\hat{\mathbf l}-\hat{\mathbf s}).
\]
Multiplying both sides by
$\hat{\mathbf j}=\hat{\mathbf l}+\hat{\mathbf s}$ and passing to the
eigenvalues gives
\[
 g_jj(j+1)=\frac12(g_l+g_s)j(j+1)
 +\frac12(g_l-g_s)[l(l+1)-s(s+1)].
\]
On setting $s=1/2$ and $j=l\pm1/2$, we obtain
\begin{equation}
 g_j=g_l\pm\frac{g_s-g_l}{2l+1}\quad\text{for }j=l\pm1/2.
 \tag{118.11}
\end{equation}
With the gyromagnetic factors quoted above, the proton magnetic moment
$\mu_p=g_jj$ is
\begin{equation}
 \begin{aligned}
  \mu_p&=\left(1-\frac{2{,}29}{j+1}\right)j
     &&\text{for }j=l-1/2,\\
  \mu_p&=j+2{,}29 &&\text{for }j=l+1/2,
 \end{aligned}
 \tag{118.12}
\end{equation}
and the neutron magnetic moment is
\begin{equation}
 \begin{aligned}
  \mu_n&=\frac{1{,}91}{j+1}j &&\text{for }j=l-1/2,\\
  \mu_n&=-1{,}91 &&\text{for }j=l+1/2.
 \end{aligned}
 \tag{118.13}
\end{equation}
(T.~Schmidt, 1937).

If only one nucleon lies outside the filled shells, (118.12) or (118.13)
directly supplies the nuclear magnetic moment. The addition of the magnetic
moments is also elementary for two nucleons (see Problem 1). With more
nucleons, the magnetic moment must be averaged using a system wave function
properly constructed from the individual nucleon wave functions. The nucleon
configuration together with the state of the whole nucleus determines this
construction uniquely whenever the configuration admits only a single system
state with the specified $J$ and $T$ (see, for example, Problem 3). Otherwise
the nuclear state is a mixture of several independent states with the same
$J$ and $T$, and the coefficients in the linear combination that forms the
nuclear wave function are generally unknown\footnote{In practice, however,
the accuracy of the ``one-particle'' scheme for calculating nuclear magnetic
moments is not high. The paired values in (118.12) and (118.13) serve more as
upper and lower bounds than as exact moment values.}.

\SourcePage{588}{591}
Finally, the presence of nucleon spin--orbit coupling in a nucleus gives the
protons an additional magnetic moment beyond (118.9) (M.~G\"oppert-Mayer,
J.~H.~Jensen, 1952). The reason is that, when an interaction operator depends
explicitly on a particle's velocity, an external field is introduced by the
replacement $\hat{\mathbf p}\to\hat{\mathbf p}-(e/c)\mathbf A$. Making this
replacement in (118.3), and using (111.7) for the vector potential, produces
the following additional term in the proton Hamiltonian:
\[
 \varphi(r)\frac{e}{cm_p}(\mathbf n\times\mathbf A)\hat{\mathbf s}
 =f(r)\frac{e}{2c\hbar}\{\mathbf r\times(\mathbf H\times\mathbf r)\}
   \hat{\mathbf s}
 =f(r)\frac{e}{2c\hbar}
   \{\mathbf r\times(\hat{\mathbf s}\times\mathbf r)\}\mathbf H.
\]
This term is equivalent to an additional magnetic moment whose operator is
\begin{equation}
 \hat{\boldsymbol\mu}_{\text{add}}
 =-\frac{e}{2c\hbar}f(r)
   \{\mathbf r\times(\hat{\mathbf s}\times\mathbf r)\}
 =-\frac{e}{2c\hbar}r^2f(r)
 \{\hat{\mathbf s}-(\hat{\mathbf s}\mathbf n)\mathbf n\}.
 \tag{118.14}
\end{equation}

\ProblemHeading 1. Find the magnetic moment of a two-nucleon system, whose
total mechanical angular momentum is $\mathbf J=\mathbf j_1+\mathbf j_2$,
expressing it in terms of the individual nucleon magnetic moments $\mu_1$ and
$\mu_2$.

\SolutionHeading A derivation analogous to that of (118.11) gives
\[
 \frac\mu J=\frac12\left(\frac{\mu_1}{j_1}+\frac{\mu_2}{j_2}\right)
 +\frac12\left(\frac{\mu_1}{j_1}-\frac{\mu_2}{j_2}\right)
 \frac{(j_1-j_2)(j_1+j_2+1)}{J(J+1)}.
\]

\ProblemHeading 2. Determine the possible states of three nucleons with
angular momenta $j=3/2$ and identical principal quantum numbers.

\SolutionHeading We proceed as in \S~67 when the possible states
of equivalent electrons were found. Each nucleon may occupy one of eight
states, characterized by the following pairs $(m_j,\tau_\zeta)$:
\[
 \begin{gathered}
 (3/2,1/2),\quad(1/2,1/2),\quad(-1/2,1/2),\quad(-3/2,1/2),\\
 (3/2,-1/2),\quad(1/2,-1/2),\quad(-1/2,-1/2),\quad(-3/2,-1/2).
 \end{gathered}
\]
Combining three distinct states at a time, we obtain the following pairs
$(M_J,T_\zeta)$ for the three-nucleon system:
\[
 (7/2,1/2),\quad2(5/2,1/2),\quad(3/2,3/2),\quad4(3/2,1/2),
 \quad(1/2,3/2),\quad5(1/2,1/2).
\]
The numeral preceding a pair gives the number of corresponding states; states
with negative $M_J$ and $T_\zeta$ need not be listed. These pairs correspond
to system states with the following $(J,T)$ values:
\[
 (7/2,1/2),\quad(5/2,1/2),\quad(3/2,3/2),\quad(3/2,1/2),\quad(1/2,1/2).
\]

\ProblemHeading 3. Taking isotopic invariance into account, determine the
magnetic moment of the ground state for a configuration consisting of two
neutrons and one proton in $p_{3/2}$ states with the same
$n$\footnote{The ${}^{7}\mathrm{Li}$ nucleus has this configuration outside
the filled shell $(1s_{1/2})^4$.}.

\SourcePage{589}{592}
\SolutionHeading The ground state of this configuration has
$J=3/2$, while the rule stated in the text gives its isospin the smallest
possible value, $T=|T_\zeta|=1/2$.

We determine the system wave function belonging to the largest possible value
$M_J=3/2$. Subject to the Pauli-principle requirements for the two identical
nucleons, this $M_J$ value may be realized by the following triples of $m_j$
values for the nucleons $p,n,n$, respectively:
\[
 (3/2,3/2,-3/2),\quad(3/2,1/2,-1/2),\quad
 (1/2,3/2,-1/2),\quad(1/2,3/2,1/2).
\]
The required wave function $\Psi^{JM_J}_{TT_\zeta}$ is therefore a linear
combination of the form
\begin{equation}
 \begin{split}
 \Psi^{3/2\;3/2}_{1/2\;-1/2}
 &=a[\psi^{3/2}_{1/2}\psi^{3/2}_{-1/2}\psi^{-3/2}_{-1/2}]
 +b[\psi^{3/2}_{1/2}\psi^{1/2}_{-1/2}\psi^{-1/2}_{-1/2}]\\
 &\quad+c[\psi^{-3/2}_{-1/2}\psi^{1/2}_{1/2}\psi^{-1/2}_{-1/2}]
 +d[\psi^{3/2}_{-1/2}\psi^{1/2}_{-1/2}\psi^{-1/2}_{1/2}],
 \end{split}
 \tag{1}
\end{equation}
where $[\ldots]$ denotes a normalized antisymmetrized product, that is, a
determinant of the form (61.5), of the individual nucleon wave functions
$\psi^{m_j}_{\tau_\zeta}$.

The function (1) must vanish under the action of each of the operators
\[
 \hat T_- =\sum_{i=1}^3\hat\tau_-^{(i)}
 \quad\text{and}\quad
 \hat J_+=\sum_{i=1}^3\hat j_+^{(i)}
\]
(compare the problem in \S~67). The operators $\hat\tau_-^{(i)}$ turn the
proton function of the $i$th nucleon into a neutron function (and annihilate a
neutron function). It is therefore readily seen that $\hat T_-$ turns the
first term in (1) into a determinant with two identical rows and hence into
zero, while the determinants from the other three terms become identical. We
thus obtain the condition $b+c+d=0$. Next, for a single nucleon with $j=3/2$
and the various $m_j$ values, (27.12) gives
\[
 \hat j_+\psi^{3/2}=0,\quad
 \hat j_+\psi^{1/2}=\sqrt3\psi^{3/2},\quad
 \hat j_+\psi^{-1/2}=2\psi^{1/2},\quad
 \hat j_+\psi^{-3/2}=\sqrt3\psi^{-1/2}.
\]
It then follows that applying $\hat J_+$ to (1) yields
\[
 \hat J_+\Psi^{3/2\;3/2}_{1/2\;-1/2}
 =\sqrt3(a+b-c)
 [\psi^{3/2}_{1/2}\psi^{3/2}_{-1/2}\psi^{-1/2}_{-1/2}]
 +2(c-d)[\psi^{3/2}_{-1/2}\psi^{1/2}_{1/2}\psi^{1/2}_{-1/2}].
\]
The sign changes in some terms result from interchanging rows of the
determinant. Requiring this expression to vanish gives
\[
 a+b-c=0,\qquad c-d=0.
\]
Together with normalization of (1), these relations yield
\[
 a=\frac3{\sqrt{15}},\qquad b=-\frac2{\sqrt{15}},\qquad
 c=d=\frac1{\sqrt{15}}.
\]
The mean value of the proton (or neutron) magnetic-moment projection in a
state with given $m_j$ is $\mu_pm_j/j$ (or $\mu_nm_j/j$). The mean system
moment calculated with (1) is consequently
\[
 \overline\mu=\overline\mu_z
 =\frac9{15}\mu_p+\frac4{15}\mu_p
 +\frac1{15}\left(\frac13\mu_p+\frac23\mu_n\right)
 +\frac1{15}\left(-\frac13\mu_p+\frac34\mu_n\right)
 =\frac1{15}(13\mu_p+\mu_n).
\]
Equations (118.12) and (118.13) give, for a nucleon in a $p_{3/2}$ state,
$\mu_n=-1{,}91$ and $\mu_p=3{,}79$. The result is $\mu=3{,}03$.
